\documentclass[11pt]{article}

% ---------------------------------------------------------------------------
%  The aligned padded frame of a Tucker tensor train and the cone reference:
%  completing frames at a rank increase by truncating A + eps S with projection-only
%  access to S, specialized to T3s -- the mathematical companion of the user guide
%  docs/aligned_padded_frame.md and the design record
%  docs/contributor/padded_frame_internals.md (2026.4.0).
%
%  Compile with pdflatex (needs: amsmath, amssymb, amsthm, mathtools, algorithm,
%  algpseudocode, booktabs, url, hyperref, geometry -- all standard):
%      latexmk -pdf padded_frame_cone.tex
%
%  It specializes the maintainer's note "Completing frames by truncating A + eps S, with
%  projection-only access" (general trees) to T3s, and cites the implemented functions.
% ---------------------------------------------------------------------------

\usepackage[margin=1in]{geometry}
\usepackage[T1]{fontenc}
\usepackage{amsmath,amssymb,amsthm}
\usepackage{mathtools}
\usepackage{algorithm}
\usepackage{algpseudocode}
\usepackage{booktabs}
\usepackage{url}
\usepackage[colorlinks=true,linkcolor=blue,citecolor=blue,urlcolor=blue]{hyperref}

\makeatletter
\def\UrlBreaks{\do\_\do\.\do\/\do\(\do\,}
\makeatother
\Urlmuskip=0mu plus 2mu
\DeclareUrlCommand\code{\urlstyle{tt}}

\newcommand{\R}{\mathbb{R}}
\newcommand{\T}{\mathsf{T}}
\newcommand{\Tc}{\mathcal{T}}
\newcommand{\Nc}{\mathcal{N}}
\newcommand{\Jac}{\mathcal{J}}
\newcommand{\ip}[2]{\langle #1,\, #2\rangle}
\newcommand{\eps}{\varepsilon}
\newcommand{\Abar}{\bar A}
\newcommand{\Impl}[1]{\Statex \hspace*{-\algorithmicindent}\parbox[t]{\dimexpr\linewidth+\algorithmicindent\relax}{\raggedright\textbf{Code:} #1}}

\theoremstyle{definition}
\newtheorem{definition}{Definition}
\newtheorem{remark}{Remark}
\theoremstyle{plain}
\newtheorem{proposition}{Proposition}

\title{The aligned padded frame of a Tucker tensor train,\\ and the cone reference}
\author{T3Toolbox 2026.4.0}
\date{}

\begin{document}
\maketitle

\begin{abstract}
At a zero-padded rank increase, a Tucker tensor train (T3) keeps its value but its orthonormal frame gains new
directions in every rank store, and nothing about the point determines them. We choose them from a \emph{hint} $S$,
a tensor accessible only through projections onto tangent spaces (in a fit, the misfit gradient). The construction
truncates $A+\eps S$ to leading order in a formal $\eps$, which fills the new directions from the normal blocks of
$S$ and keeps $A$'s frames exactly; with projection-only access it becomes a fixed-point iteration whose step is
unrooted message passing over the T3's directed edges, and whose factorizations are no wider than the old rank plus
twice the increase. We state what its fixed points satisfy, why they are local and not global, and how the library
implements the start, the step, the schedule and the cone reference $\|\Pi_{T(F)} \nabla f(0)\|$ of the Newton stop.
\end{abstract}

\tableofcontents

% ===========================================================================
\section{The problem}\label{sec:problem}

Rank continuation fits a T3 $X_{\mathrm{old}}$ at some ranks, raises the ranks, and starts Newton at
$X_0 = X_{\mathrm{old}}$ zero-padded. Newton linearizes at a frame of $X_0$: orthonormal frame cores in four
\emph{stores} (up, down, left, right), whose tangent space $T(F)$ holds the first step. The old frame determines
the leading slots of every store; the new slots are a free orthonormal completion, and the first step from $X_0$
depends on it. A completion that ignores the data can leave the step blind to the directions the data want.

Given a hint $S$ (a tensor; in a fit, the misfit gradient $\nabla f(X_0)$), we want completions whose tangent space
holds much of $S$. The construction is that of the maintainer's note ``Completing frames by truncating $A+\eps S$, with
projection-only access'' \cite{AlgerNote} (\emph{the note} below), written for general tree tensor networks; this
document specializes it to T3s and to the library's implementation. The T3 frames, tangent spaces and variations are
those of T4S \cite{T4S}. Two constraints shape the answer: $S$ is available only through projections
$F\mapsto\Pi_{T(F)}S$, at the cost of one gradient each; and the old frames must survive bitwise, so that the frame
still represents $X_0$.

\paragraph{Notation.} $d$ modes $i=0,\dots,d-1$ of sizes $N_i$. $A=X_{\mathrm{old}}$, the old point. $T$ is the
tangent space of the fixed-rank manifold at $A$, and $N=\Pi_{T^\perp}S$ the normal part of $S$. For a frame $F$ at
the new ranks, $T(F)$ is its tangent space. $\|\cdot\|$ is the Frobenius norm.

% ===========================================================================
\section{Frames of a T3, and padded frames}\label{sec:frames}

A T3 is a tree: Tucker leaves $\ell_i$ (physical size $N_i$) and TT cores $c_i$ in a chain (the orthogonal frame of a
T3 is T4S \cite{T4S} App.~A.1, A.3, Algorithm~11). For a directed edge
$x\to t$, the frame $U_{x\to t}$ is an orthonormal basis of the column space of the tensor unfolded with $x$'s side
as rows. Frames are nested: $U_{x\to t}$ lies in $H_x\otimes\bigotimes_{y\in\mathrm{nb}(x)\setminus t}U_{y\to x}$,
with $H_x$ the physical space at $x$ (trivial at a TT core). So one \emph{frame core} per directed edge stores the
frame, and the four directions are exactly the library's four stores (Table~\ref{tab:stores}). The stores' widths
are \code{ranks.compute_orthogonal_representation_ranks} of the ranks.

\begin{table}[h]
\centering
\begin{tabular}{@{}llll@{}}
\toprule
directed edge & store & core & output slot \\
\midrule
$\ell_i\to c_i$ & \code{up[i]} & $U_i\in\R^{n_{U,i}\times N_i}$ & the Tucker index \\
$c_i\to \ell_i$ & \code{down[i]} & $O_i\in\R^{r_{L,i}\times n_{D,i}\times r_{R,i+1}}$ & the down index \\
$c_i\to c_{i+1}$ & \code{left[i]}, $i\le d-2$ & $L_i\in\R^{r_{L,i}\times n_{U,i}\times r_{L,i+1}}$ & bond $i+1$ \\
$c_i\to c_{i-1}$ & \code{right[i]}, $i\ge 1$ & $R_i\in\R^{r_{R,i}\times n_{U,i}\times r_{R,i+1}}$ & bond $i$ \\
\bottomrule
\end{tabular}
\caption{Directed edges of a T3 and the frame stores. $L_{d-1}$ and $R_0$ are the \emph{carriers} of the base
point, not frame cores of an edge.}\label{tab:stores}
\end{table}

\begin{definition}[padded frame]
A \emph{padded frame} of $A$ at larger stores has, in every store, the old frame core bitwise in the leading slots
(the \emph{A-block}, $\Abar_{x\to t}$, of width $r_e$) and an orthonormal \emph{completion} $P_{x\to t}$ of width
$k_e$ in the trailing slots (the \emph{P-block}). An A-block core is zero unless every input slot is in its
A-part. The carriers hold $A$'s content only.
\end{definition}

Every padded frame represents $X_0$, whatever the completion: the two base-point chains run through A-parts only.
The \emph{empty frame} (\code{padded_frames.fv_empty_frame}, \code{uniform_padded_frames.ufv_empty_frame}) is the
orthogonal representation of the rank-0 T3, Tucker ranks $(0,\dots,0)$ and TT ranks $(1,0,\dots,0,1)$: zero-width
rank axes except the boundary bonds. It is the old frame of the zero tensor, so a start or an alignment at the zero
tensor is the general construction with no A-blocks, not a special case.

% ===========================================================================
\section{Completing frames by truncating \texorpdfstring{$A+\eps S$}{A + eps S}}\label{sec:truncation}

Truncating $A$ to larger ranks completes its frames arbitrarily. Truncate $A+\eps S$ instead, with $\eps$ formal and
the two parts stored separately, and keep leading orders: $S$ fills the missing directions at order $\eps$, and $A$'s
frames survive exactly.

\paragraph{Leading orders, and the tilt.} Dual numbers ($\eps^2=0$) do not suffice: the new basis vectors come out as
$\eps v$ and cannot be normalized. Leading-order arithmetic does (write each quantity as $\eps^p(c+O(\eps))$, compare by
$p$ and then $c$). One trap remains. For $A=U\Sigma V^\T$, the leading frame of $A+\eps S$ is
$U+\eps(I-UU^\T)SV\Sigma^{-1}+O(\eps^2)$: the tangent part of $S$ \emph{tilts} $A$'s frame, and that tilt is not a
new direction. The new directions must come from the normal block, with $A$'s frames removed on both sides,
\begin{equation}\label{eq:fill}
  \text{fill with } (I-UU^\T)\,S\,(I-VV^\T), \qquad\text{not with } (I-UU^\T)\,S .
\end{equation}
$\Sigma^{-1}$ cancels from the fill, so $A$'s singular values never enter and ill-conditioned edges of $A$ play no
role. In the matrix case the left completion is the top $k$ left singular vectors of the normal part $N$, and the
right one its top $k$ right singular vectors: the truncated SVD of $N$, which is optimal.

\paragraph{The sweep on a tree.} In a leaves-to-root sweep where each edge sees the earlier truncations, every SVD
step of the formal truncation reduces to an SVD of one normal block of $S$ with its rows restricted to the candidate
space: at edge $x\to t$, the rows in $C'_{x\to t}\ominus U_{x\to t}$ with $C'$ built from the completed child frames,
and the columns in the complement of $A$'s frame $U_{t\to x}$. Its completions depend on the root: each edge sees the
choices made below it. On small test problems the five roots of a five-site TT captured between $0.47$ and $0.54$ of
$\|N\|^2$, against $0.59$ for the best completion found. The loop of the next section removes that dependence.

% ===========================================================================
\section{Projection-only access: the alignment loop}\label{sec:loop}

When $S$ can only be projected, run the formal truncation on $A+\eps Z$, where $Z=\Pi_{T(F)}S$ is the projection of
$S$ onto the tangent space of the current padded frame, and repeat:
\begin{equation}\label{eq:loop}
  F \;\leftarrow\; \operatorname{step}\bigl(F,\;\Pi_{T(F)}S\bigr).
\end{equation}
Each iteration touches $S$ once; the step works on $Z$, a T3 tangent vector, through its gauged variation cores (T4S
eqs.~(46)--(48)), never forming it as the doubled-rank T3 it is (T4S App.~A.3.1, eqs.~(49)--(52)).

\paragraph{Why the two-sided tangent space.} Projecting $S$ onto the linear piece a rooted completion defines (write
$A$ in the completed frames, vary one core at a time) returns the same completions: every frame set is a fixed point,
and the loop stalls (measured: angles of $10^{-15}$). The tangent space $T(F)$ of the padded point also holds, at each
edge, tensors whose far side lies in the \emph{far completion} $Q_e$ (the completion of $U_{t\to x}$). Through them
the normal block of $S$ applied to $Q_e$ enters the step: information from outside the current completions.

\paragraph{The matrix case.} With $P,Q$ completing $U,V$, the projection of $S$ onto $T(F)=\{U'X+YV'^\T\}$ has normal
block $PP^\T N+(I-PP^\T)NQQ^\T$. The new $P$ is its top-$k$ left singular subspace, and the new $Q$ the top-$k$ right
singular subspace of the captured part. This is a two-sided subspace iteration on $N$; its only stable fixed point is
the top-$k$ singular pair, and the error contracts by about $\sigma_{k+1}/\sigma_k$ per iteration.

% ===========================================================================
\section{What the fixed points satisfy}\label{sec:fixed}

\begin{definition}[fixed point]
$F$ is a fixed point of \eqref{eq:loop} when the step returns the same completions: at every directed edge, the
completion spans the leading new directions of its own content, given the other edges' completions.
\end{definition}

Equivalently (the note's stopping criterion), at every edge the content entering through the far completion already
lies in the near completion, provided the $k\times k$ matrix $P_{a\to b}^\T S P_{b\to a}$ is invertible. In the
library's terms: one more pass moves nothing, and a pass updating any one store alone (the others transported,
\S\ref{sec:schedule}) moves nothing. The docstring of \code{fv_aligned_padded_frame} demonstrates both.

\subsection{The captured energy: stationary, not maximal}

\begin{proposition}[matrix case]\label{prop:matrix}
Let $N=(I-UU^\T)S(I-VV^\T)$, and let $P\perp U$, $Q\perp V$ be completions of width $k$. Then
\begin{equation}\label{eq:energy}
  E(P,Q)\;:=\;\|\Pi_{T(F)}S\|^2 \;=\; \|S\|^2-\bigl\|(I-PP^\T)\,N\,(I-QQ^\T)\bigr\|^2 .
\end{equation}
If $\sigma_k>\sigma_{k+1}>0$, the fixed point $(P,Q)=$ the top-$k$ singular pair of $N$ is a stationary point of $E$
and a saddle. Rotating $P$ alone (or $Q$ alone) out of its span never decreases $E$, and increases it at second order
when the rotation has a component along some $u_j$ (or $v_j$), $j>k$; rotating both together toward the next singular
pair decreases it at second order.
\end{proposition}

\begin{proof}
$\Pi_{T(F)}=I-(I-U'U'^\T)(\cdot)(I-V'V'^\T)$ with $U'=[U\;P]$, $V'=[V\;Q]$, and
$(I-U'U'^\T)S(I-V'V'^\T)=(I-PP^\T)N(I-QQ^\T)$ since $P\perp U$, $Q\perp V$; this gives \eqref{eq:energy}. With $Q$
fixed at the top-$k$ right singular vectors, $M:=N(I-QQ^\T)=\sum_{j>k}\sigma_ju_jv_j^\T$ and
$E=\|S\|^2-\|M\|^2+\|P^\T M\|^2$. At $P=[u_1\cdots u_k]$, $P^\T M=0$, the minimum of $\|P^\T M\|^2\ge 0$: so $E$ is
stationary in $P$ and does not decrease under any rotation of $P$ alone. A rotation of $u_k$ by $h$ toward a unit
vector $w\perp[U\;P]$ raises $\|P^\T M\|^2$ by $h^2\|w^\T M\|^2+O(h^3)$, which is positive when $w$ has a component
along some $u_j$, $j>k$ (a $w$ in the left null space of $N$ leaves $E$ unchanged). Symmetrically in $Q$. For the
paired rotation $u_k\to\cos\theta\,u_k+\sin\theta\,u_{k+1}$, $v_k\to\cos\theta\,v_k+\sin\theta\,v_{k+1}$, only the
$2\times2$ block on $(u_k,u_{k+1})\times(v_k,v_{k+1})$ changes; there the residual is the outer product of the two
complements, with coefficient $\sigma_k\sin^2\theta+\sigma_{k+1}\cos^2\theta$, so
\begin{align*}
  R(\theta):=\bigl\|(I-PP^\T)N(I-QQ^\T)\bigr\|^2
  &=\sum_{j>k+1}\sigma_j^2+\bigl(\sigma_{k+1}+(\sigma_k-\sigma_{k+1})\sin^2\theta\bigr)^2\\
  &= R(0)+2\sigma_{k+1}(\sigma_k-\sigma_{k+1})\,\theta^2+O(\theta^4).
\end{align*}
$R$ grows at second order, so $E$ falls, exactly when $\sigma_k>\sigma_{k+1}>0$; with $\sigma_{k+1}=0$ the decrease is
only fourth order.
\end{proof}

$E$ counts content of rank up to $2k$ (its maximum puts $P$ and $Q$ on \emph{different} singular directions), so it
is the wrong objective for a completion; the loop does not maximize it. What carries over to trees is stationarity
(the \emph{Galerkin-type condition}): the completions are a stationary point of $E$, so rotating one store's
completion out of its span changes the captured energy only at second order. Measured on a $d=4$ T3 at a converged
fixed point (\code{fv_aligned_padded_frame} at $\mathrm{tol}=10^{-12}$): rotating any one of ten completed stores by
$h=10^{-3}$ and $10^{-4}$ changed $E$ in the ratio $100.0$ (second order), and always upward, as in the matrix case.

\subsection{The cone merit: local maxima}

What the fixed points do maximize locally, in the measurements, is the note's \emph{merit}. For completions toward a root $c$, let $T_A(F)$ be the linear
space through $A$ obtained by writing $A$ in the completed child frames and the old parent frame and varying one core
at a time, and $\Nc_c(F)=T_A(F)\ominus T$. The merit is
\begin{equation}\label{eq:merit}
  m_c(F)=\frac{\|\Pi_{\Nc_c(F)}N\|^2}{\|N\|^2},
\end{equation}
the share of the normal part that the new directions represent on one side of every edge. Its best value is the
same for every root, because the pieces $T_A(F)$ together fill the tangent cone at $A$. In the matrix case
$m=\|P^\T N\|^2/\|N\|^2$, maximized by the top-$k$ singular subspace, with every other stationary point a saddle.

On trees the fixed points of the loop are the stationary points of $m_c$ (the information entering through the far
completions lies in the near ones, which is also the condition for $\nabla m_c=0$), so they do not depend on how the
truncation inside the step is organized. Measured on small trees: gradients below $10^{-11}$ at every fixed point;
from the same start, the unrooted step and steps rooted at different sites converged to the same frames to
$10^{-15}$; the merit never decreased along the way.

\paragraph{Summary.} At a fixed point, every store's completion spans the leading new directions of its own content,
given the other stores' completions, so one more pass moves nothing. The completions are then a stationary point of
the captured hint energy $\|\Pi_{T(F)}\,\mathrm{hint}\|^2$ (proved for matrices, measured on T3s) but generally not a
maximum of it: rotating one store's completion alone never lowers that energy and generically raises it at second
order. In measurements on small trees, they are a strict local maximum of the cone merit, one of many.

\paragraph{Many strict local maxima.} Every stationary point the loop reached was a strict local maximum of the merit
(negative definite Hessians on a slice transversal to the rotation symmetries), and there are many. Each edge spends
$k$ directions, and the excursions of $N$ compete for them. An excursion across $L$ edges is captured only if all $L$
completions align with it, so its captured energy is a product of $L$ alignment factors, of order $2L$ in the
rotation angles. Near a configuration that ignores it, turning one edge toward it loses energy at second order and
gains only at order $2L$: that configuration is a strict local maximum, however large the ignored excursion. The
minimal example (three sites, $r=k=1$, short excursions of energy $\alpha^2$ and $\gamma^2$, a long one $\beta^2$):
\begin{equation}
  m(\theta,\varphi)=\alpha^2\cos^2\theta+\gamma^2\cos^2\varphi+\beta^2\sin^2\theta\,\sin^2\varphi ,
\end{equation}
where ``both short'' ($\theta=\varphi=0$) is always a local maximum, and ``long'' ($\theta=\varphi=\pi/2$) is one
exactly when $\beta^2>\max(\alpha^2,\gamma^2)$. No local method leaves a wrong maximum, so the start decides.
Oversampled completions see the energy of long excursions through each edge, which a random start does not; this is
why the schedule oversamples first (\S\ref{sec:schedule}). Matrices have no spurious maxima (one edge).

% ===========================================================================
\section{One step on cores}\label{sec:step}

The step replaces every completion by the leading new directions of the tangent's normal content at that edge, with
the A-block pinned. It needs no factorization wider than the old width plus $2k$.

\paragraph{Messages.} On each directed edge $a\to b$, with $Z=Z_{\mathrm{near}}+Z_{\mathrm{far}}$ split into the
node terms on $a$'s side and on $b$'s side, $P$ the old completion and $P^+$ the new one:
\begin{center}
\begin{tabular}{@{}llll@{}}
\toprule
message & size & definition & built \\
\midrule
$\Phi_{a\to b}$ & $k\times R$ & $P_{a\to b}^\T Z_{\mathrm{far}}U'_{b\to a}$ & stage 1, far leaves inward\\
$\Xi_{a\to b}$ & $k\times\cdot$ & $\Xi\Xi^\T=P_{a\to b}^\T Z_{\mathrm{far}}(I-U'_{b\to a}U'^\T_{b\to a})Z_{\mathrm{far}}^\T P_{a\to b}$ & stage 1, with $\Phi$\\
$D_{a\to b}$ & $R\times k$ & $U'^{+\T}_{a\to b}Z_{\mathrm{near}}P_{b\to a}$ & stage 2, message order\\
$o_{a\to b}$ & $k^+\times k$ & $P^{+\T}_{a\to b}P_{a\to b}$ & stage 2, with $D$\\
\bottomrule
\end{tabular}
\end{center}
$\Phi$ and $\Xi$ say what the far side sends through the old completion: the part inside the old far frame, and a
square root of the rest (any factor; only $\Xi\Xi^\T$ enters). $D$ and $o$ carry the near side, already rotated into
the new frames.

\paragraph{The edge identity.} At $a\to b$, with $a$'s incoming slots in their new coordinates (the old block kept,
the completion block multiplied by $o$), let $g_{\mathrm{old}}$ be the old completion and $g_{\mathrm{new}}$ the near
content on the far completion. Then the edge's normal content factors as
\begin{equation}\label{eq:edge}
  Y=\bigl[\,g_{\mathrm{old}}\Phi_P+g_{\mathrm{new}}\;\big|\;g_{\mathrm{old}}\Xi\,\bigr]\,\Theta^\T ,
\end{equation}
with $\Theta$ orthonormal, so its left singular subspace is that of the bracket, of width at most $2k$ beside the
A-block. The new completion is the leading $k^+$ left singular directions of $(I-\Abar\Abar^\T)Y$ completed against
$\Abar$, and the edge's \emph{movement} is $\|g_{\mathrm{old}}-P^+P^{+\T}g_{\mathrm{old}}\|$: zero exactly when the
completion did not turn.

\begin{algorithm}[t]
\caption{One alignment step (ragged)}\label{alg:step}
\begin{algorithmic}[1]
\Require padded frame $F$; variations $(V,H)$ of $Z$ at $F$; old ranks; output ranks; \code{update}
\State Stage 1 (old frames): down messages; left messages from the far end; right messages from the near end; up
messages; each $(\Phi,\Xi)$.
\Impl{\code{_down_message}, \code{_left_message_step}, \code{_right_message_step} (scan bodies), \code{_up_message};
$\Xi$ by \code{_root}, a factor from the thin SVD, never a Gram matrix.}
\State Stage 2 (message order): up edges (one per sharing group, the members' contents side by side); the left sweep
$i=0..d-2$; the right sweep $i=d-1..1$; down edges. At each: new coordinates, $Y$ by \eqref{eq:edge}, $P^+$, $o$, $D$,
movement.
\Impl{\code{_new_coordinates}, \code{_edge_content}, \code{_edge} $\to$ \code{linalg.pinned_leading_columns};
\code{_up_edge}, \code{_left_edge_step}, \code{_right_edge_step}, \code{_down_edge}.}
\State A store not in \code{update} is transported: $Y=g_{\mathrm{old}}$.
\State Carriers: the old content, embedded. Return the frame and the largest movement.
\Impl{\code{padded_frames.tv_align_padded_frame}; frontend one pass: \code{T3Tangent.align_padded}.}
\end{algorithmic}
\end{algorithm}

The messages past the chain's ends are zero-size arrays, so every core runs the same body and both sweeps are scans.
A pass that updates some stores and transports others may not also truncate a transported store: a transported
completion can lose directions arbitrarily in a truncated neighbour's coordinates, so the step raises.

\paragraph{The uniform twin.} \code{uniform_padded_frames.utv_align_padded_frame} computes the same step on padded
supercores in block layout ($[\text{A-block}\mid\text{P-block}]$ per slot, the split $(N_a,R_b)$ static). The edge is
one \code{linalg.pad_safe_svd} of $[\Abar\mid Y_\perp/(4\|Y_\perp\|)]$: the pins have singular value $1$ and the scaled
content at most $1/4$, so the pins lead, and a host one-hot selects the $k^+$ triplets after them. Every count comes
from the masks; no rank is decided from content.

% ===========================================================================
\section{The start}\label{sec:start}

The loop starts from the hint-free padded frame. Every store's new slot $s$ is seeded by row $s$ of a keyed Gaussian
stream, \code{RandomState([seed, store tag, key])} (tags up 0, down 1, left 2, right 3; the key is the mode, or a
sharing group's first mode), with the ragged unfolding's row length; a start at old width $w$ uses rows
$w,\dots,w+k-1$. The seeds are completed against the A-block in Gram--Schmidt order by the sign-fixed Householder
QR of $[\Abar\mid\text{seeds}]$.

\begin{algorithm}[t]
\caption{The hint-free padded frame}\label{alg:start}
\begin{algorithmic}[1]
\Require old frame (or the empty frame); target ranks; \code{sharing}; \code{seed}
\For{each store and core}
  \State embed the old core in the leading slots; seeds $\gets$ rows $w..w+k-1$ of its keyed stream
  \State completion $\gets$ \code{pinned_completion}$(\Abar, \text{seeds}, k)$
\EndFor
\Impl{\code{padded_frames.fv_padded_frame} (\code{_seed_rows}, \code{_padded_core}, \code{_padded_tucker_factor});
uniform: \code{uniform_padded_frames.ufv_padded_frame}, the same rows placed at the real rows, completed by
\code{linalg.pad_safe_pinned_completion}; frontend \code{T3Frame.padded}, \code{UT3Frame.padded}.}
\end{algorithmic}
\end{algorithm}

The completion is a uniformly random subspace of the A-block's complement in distribution, the natural
no-information choice, and its pivots are bounded away from zero with high probability. Because the seeds are
slot-indexed, a start built from a start never reuses a seed.

\paragraph{Representation dependence.} A TT or down seed row is read in the slot coordinates of the neighbours'
frames, and a Gaussian row does not flip when a slot's sign does. So the start is a function of the old frame
\emph{as represented}: bases chosen differently at singular-value ties, or a flipped slot sign (free in an
unnormalized SVD), give different starts, all valid. The ragged and uniform orthogonalizations of one point choose
differently, and their starts differ (span distances $0.95$--$1.00$ measured, at generic, \code{t3svd} and shared
points), as can starts on different BLAS builds. Tucker completions are seeded in physical rows and are unaffected.
The aligned output does not have this dependence, up to its tolerance (\S\ref{sec:measured}).

% ===========================================================================
\section{The schedule}\label{sec:schedule}

\begin{algorithm}[t]
\caption{The aligned padded frame}\label{alg:loop}
\begin{algorithmic}[1]
\Require old frame; target ranks; hint $S$ and \code{project}; $p$, $n_{\mathrm{over}}$, $\mathrm{cap}$, tol,
\code{alternate}
\State working $\gets$ \code{oversampled_ranks}(target, $p$); $n_{\mathrm{over}}\gets 0$ if the working stores equal
the target stores; $n_{\mathrm{over}}\gets\min(n_{\mathrm{over}},\mathrm{cap}-1)$
\State $F\gets$ the start (Algorithm~\ref{alg:start}) at the working ranks; \textbf{if} $\mathrm{cap}=0$ \textbf{return} it at the target
\For{$n_{\mathrm{over}}$ passes} $F\gets\operatorname{step}(F,\Pi_{T(F)}S)$ at the working ranks, all stores \EndFor
\State $F\gets\operatorname{step}(F,\Pi_{T(F)}S)$, output at the target ranks (truncating), all stores
\While{a whole pair fits under the cap}
  \State $F\gets$ step updating up and left (\code{'UL'}); $F\gets$ step updating right and down (\code{'RO'})
  \State \textbf{if} both movements $<$ tol \textbf{then stop}
\EndWhile
\State \textbf{if} one slot is left and not stopped: one more \code{'UL'} step
\Impl{\code{padded_frames.fv_aligned_padded_frame} (host loop; \code{use_jit}: a memoized jitted step);
\code{uniform_padded_frames.ufv_aligned_padded_frame} (\code{use_jit}: the whole loop as one \code{jax.jit},
compacted to the level's widths at exit); frontend \code{frame_variations_format.fv_aligned_padded_frame},
\code{uniform_frame_variations_format.ufv_aligned_padded_frame}.}
\end{algorithmic}
\end{algorithm}

\paragraph{Oversample, then truncate.} Completions of width $k+p$ keep enough of $S$ that the loop settles near what a
single sweep on $S$ itself would give, in a good basin (\S\ref{sec:fixed}); the truncating pass then drops to width $k$
(the same step with fewer selected directions), and the passes after it climb to that basin's fixed point. On a toy
family, $p=5$ for two passes reached $97\%$ of the converged merit in four passes and $99.7\%$ in six; oversampling
alone plateaus at $88\%$.

\paragraph{Two colours.} The simultaneous update (all stores per pass) can cycle with period~$2$: measured on a Darcy
gradient from about pass $400$, and on a $243$-number case whose stationary point lies above the cycle and is not an
attractor of that map. Updating the stores pointing toward the root $c_{d-1}$ (up, left) and then those pointing away
(right, down), the other colour transported, is an alternating (Gauss--Seidel) update with a projection between the
halves. Each updated store maximizes the captured energy $\|P^{+\T}Y\|^2$ of its own content $Y$ given the others;
since $Y$ depends on the store's previous completion, this is a self-consistency condition, not ascent of one
objective, and in particular not of $E$, which is a saddle at the fixed point (Proposition~\ref{prop:matrix}). The
alternating update converged on the cycling case with no merit decrease. The stop is tested on pairs only, never on the truncating pass,
and a NaN movement never stops the loop.

\paragraph{Defaults.} $p=5$ for $2$ passes, cap $10$, tol $0.05$: two oversampled passes, the truncating pass, three
pairs and a final \code{'UL'}. The movement tracks the true frame movement within $15\%$ from $1$ to $10^{-15}$. A pass
costs at most about $10\%$ of a gradient.

% ===========================================================================
\section{The cone reference}\label{sec:cone}

\code{newton_cg}'s relative stop $\|g\|\le\mathrm{gtol\_rel}\,\|g_0\|$ needs a reference scale. With
\code{g0_reference='cone'} it is the frame-gradient norm at the zero tensor of $x_0$'s ranks, on the frame aligned
with the data gradient there:
\begin{equation}
  \|g_0\| = \bigl\|\Pi_{T(F_0)}\nabla f(0)\bigr\|,\qquad F_0=\text{Algorithm~\ref{alg:loop} from the empty frame,
  hint } \nabla f(0)=\Jac^\T r(0).
\end{equation}
It is like-for-like with the gradient norms the stop compares it to, depends only on $x_0$'s ranks, and excludes the
regularizer (as the default \code{'zero'} reference does). The loop's frame approximates a local maximum of the merit, so the
value is a data-aware scale, not an intrinsic constant. With zero passes it is evaluated at the zero tensor's
hint-free start, a different frame from \code{'zero'}'s (\code{fv_zero_point_frame}), so the two values differ.

\begin{algorithm}[h]
\caption{The cone reference, and a problem's aligned start}\label{alg:cone}
\begin{algorithmic}[1]
\State $r_0\gets$ forward$(0)-$data
\State $F_0\gets$ the geometry's \code{aligned_padded_frame} at $x_0$ from the zero tensor, with \code{project} $=$
\code{_misfit_gradient_at} and hint $(\text{problem},r_0)$
\State $g\gets$ \code{_misfit_gradient_at} at $F_0$; \Return $\ip{g}{g}_{F_0}^{1/2}$
\Impl{\code{backend.optimizers._cone_reference}; the start after a rank increase (\code{padded_from=x_old}) is the
same member call from $x_{\mathrm{old}}$ with the residual at $x_0$: \code{backend.optimizers.padded_start}.}
\end{algorithmic}
\end{algorithm}

% ===========================================================================
\section{Numerical notes}\label{sec:numerics}

\begin{itemize}
\item $A$ enters only through its orthonormal frame cores; its singular values never appear.
\item The square-root form: $\Xi$ is a factor of the far remainder itself, and $Y$ is formed directly, so no Gram
matrix is squared. Gram-matrix messages lose $u\,s^2$ when the hint's tangent part is $s$ times its normal part; the
square-root form follows $u\max(1,s)$, the floor set by rounding $S$ itself (prototype, emulated low precision).
\item \code{pinned_leading_columns} projects $Y$ off $\Abar$ once and takes the thin SVD; its pinned completion's QR of
$[\Abar\mid U_k]$ then makes $P^+$ orthogonal to $\Abar$ to roundoff.
\item No rank cutoff anywhere: when $Y$ has fewer than $k^+$ nonzero singular values the trailing directions are a
valid, arbitrary completion (roundoff- and backend-dependent); they carry no content and are dropped at truncation.
\item The step is invariant to the hint's scale ($10^{-14}$ to $10^{14}$ tested) and to the variations' gauge.
\end{itemize}

% ===========================================================================
\section{What is measured, and what is open}\label{sec:measured}

\paragraph{Measured.} The step agrees with the dense per-edge update to $3\cdot10^{-14}$ (prototype) and within the
tests' $10^{-9}$ bar at $d\in\{1,2,3,5\}$, with sharing, truncation, stacks and the empty frame; the uniform twins equal
the ragged ones per element to rounding. From starts differing by $0.95$--$1.00$ in span, the two layers' aligned frames
agreed to $7\cdot10^{-4}$--$3\cdot10^{-2}$ at the default tol (equal pass counts) and to $8\cdot10^{-13}$--$6\cdot10^{-11}$
at tol $10^{-10}$ ($d=4$, nine cases). On a Darcy rank-continuation ladder the shipped alignment captured
$3$--$8\times$ the frame-gradient norm of the hint-free frame (levels $1\to6$, $6\to11$, $11\to16$). Through full
ladders with the exact Hessian (A100, 2026-11-10), the \code{'cone'} reference roughly halved the Newton count against
\code{'zero'} at the same relative tolerance (a looser stop, test errors within $4\%$), and within \code{'cone'} the
aligned start took about $30\%$ fewer CG iterations than the library's and the hint-free starts on 10 points $\times$ 100
directions, and 166 Newton / 4130 CG against the hint-free start's 191 / 4719 on 20 $\times$ 1000, with identical test
error. Every alignment stopped on the 10-pass cap, never on tol $0.05$.

\paragraph{Open.}
\begin{itemize}
\item Proofs for trees: the match between fixed points and stationary points of the merit, and the merit's
monotonicity, are measured, not proved.
\item A guarantee that the oversampled phase lands in the best basin, and a rule for $p$ and the switch tied to the
gaps of the normal blocks' spectra.
\item Flat normal spectra converge slowly (the matrix rate is $\sigma_{k+1}/\sigma_k$); near-tied hint spectra at $d=24$
did not converge in ten passes.
\item Moves between basins: a discrete move rotating one direction on every edge of an excursion at once.
\item Canonical bases for the representation, which would make the hint-free start representation-independent.
\end{itemize}

\begin{thebibliography}{9}
\bibitem{AlgerNote} N.~Alger, \emph{Completing frames by truncating $A+\varepsilon S$, with projection-only access},
  unpublished note, 2026.
\bibitem{T4S} N.~Alger, B.~Christierson, P.~Chen, O.~Ghattas, \emph{Tucker Tensor Train Taylor Series},
  arXiv preprint arXiv:2603.21141v2, 2026. \url{https://arxiv.org/abs/2603.21141v2}
\end{thebibliography}

\end{document}
